\section{Semantic Analysis and Validation}\label{sec:semantic-rules}

After the frontend produces a syntactically correct Abstract Syntax Tree, the compiler enters the \textit{Semantic
Analysis} phase.
This stage is responsible for ensuring that the program is logically and musically sound.
It performs three primary tasks: identifier resolution (scoping), \textit{explicit type checking} with expression
type propagation, and musical constraint validation.
Unlike version~1.0, the analyzer does not infer variable types from untyped \texttt{let} bindings; every variable and
pattern parameter carries a declared \texttt{types::Type}, and mismatches are reported at the binding or call site.

\subsection{Single-Pass Traversal Architecture}\label{detail-subsec:single-pass-traversal}

The semantic analyzer is implemented in the compiler's \texttt{Traversal} class.
Because the AST uses \texttt{std::variant} for statement and expression polymorphism, the analyzer relies on a visitor
pattern implemented with \texttt{std::visit} and a small overloaded-helper template.
This allows the compiler to dispatch logic cleanly based on the concrete type of the AST node without relying on dynamic
casting or virtual method tables.

\texttt{Traversal::run()} executes one coherent walk over the program:

\noindent \textbf{Header validation.} Tempo and time-signature literals are range-checked before any track is visited.

\noindent \textbf{Global scope.} Top-level patterns are registered first (\texttt{visit\_globals}), then global
statements and tracks are analyzed.
Pattern collection at scope entry is not a separate compilation pass; it is the opening step of the same traversal
that later visits bodies and expressions.

\noindent \textbf{Track and voice entry.} For each \texttt{TrackDefinition}, the analyzer enters a new lexical scope,
collects track-local patterns, then traverses track items.
The same structure applies to \texttt{VoiceDefinition} nodes.

\noindent \textbf{Bottom-up expressions.} A binary expression node recursively visits its children, obtains their
types, validates operand compatibility, and stores the computed result type in \texttt{Annotations}.

Version~1.0 relied on an additional pass (\texttt{validate\_pattern\_instantiation}) that cloned scopes and re-walked
pattern bodies at each call.
Version~2.0 removes that pass: pattern calls are validated when \texttt{visit\_call} evaluates argument expressions and
invokes \texttt{validate\_call} against the overload selected by type signature.

\subsection{Symbol Table and Scope Management}\label{detail-subsec:symbol-table-and-scope-management}

Scoping is managed by the \texttt{SymbolTable} and \texttt{ScopeStack} classes.
The \texttt{SymbolTable} maintains the actual symbol data, while the \texttt{ScopeStack} manages the hierarchical
environment during the AST traversal.

\subsubsection{Hierarchical Scoping Rules}

The \texttt{ScopeStack} maintains a stack of \texttt{ScopeId} references pointing to the \texttt{SymbolTable}.
Scopes are pushed and popped as the visitor enters and exits structural nodes (global, track, voice, pattern, and
control-flow blocks).
The implementation enforces strict visibility rules:

\noindent \textbf{Lookup.} When an \texttt{IdentifierExpression} is encountered, the \texttt{ScopeStack} performs a
reverse search from the current scope up to the global scope.
If the identifier is found, its \texttt{SymbolId} and the \texttt{types::Type} stored on the symbol are retrieved.

\noindent \textbf{Symbol metadata.} Every registered symbol stores its \texttt{SymbolKind} (variable, parameter,
pattern, track, or voice), its declared \texttt{types::Type}, its \texttt{Location} in the source file, and a raw
pointer back to its declaration node in the AST\@.

\subsubsection{AST Annotation}

To separate analysis from execution, the semantic pass does not mutate the original AST structure.
Instead, it populates an \texttt{Annotations} object.
Upon successful resolution of an identifier, \texttt{Annotations} maps the memory address of the
\texttt{IdentifierExpression} to its resolved \texttt{SymbolId}.
Every evaluated expression is mapped to its computed \texttt{types::Type}.
The lowering pass reads these cached types and symbol links; it does not repeat string lookups or re-derive types.

\subsection{Formal Type Rules and Constraints}\label{detail-subsec:formal-type-rules-and-constraints}

The type system acts as the musical guardrail of the compiler, implemented in \texttt{type\_rules.hpp} and
\texttt{type\_rules.cpp}.
All types are represented by the \texttt{motivo::types::Type} enumeration (not the legacy \texttt{TypeKind} name from
version~1.0).

\subsubsection{Type Categories}

\noindent \textbf{Scalars.} \texttt{Int}, \texttt{Double}, \texttt{Bool} --- used for loop bounds, arithmetic, and
conditionals.
Only \texttt{int}, \texttt{double}, and \texttt{bool} may appear in \texttt{VarDeclStatement} and pattern parameter
positions.

\noindent \textbf{Musical primitives.} \texttt{Note}, \texttt{Rest}, \texttt{Chord}, \texttt{Sequence}, \texttt{Drum}
--- the playable or composable entities that lowering eventually flattens into \texttt{NoteEvent} records.
Declarable variable types include \texttt{note}, \texttt{seq}, and \texttt{chord}; \texttt{rest} and drum literals are
expression-only.

\noindent \textbf{Utility types.} \texttt{Unknown} (error recovery) and \texttt{Void} (non-value statements).

\subsubsection{Strict Declarations and Assignments}

The design deliberately rejects implicit widening at declaration and assignment sites.
The function \texttt{is\_assignable(target, source)} returns true only when \texttt{target == source}.
Thus \texttt{int x = 1.0;} is rejected even though arithmetic expressions may promote \texttt{Int} and
\texttt{Double} operands via \texttt{numeric\_result}.
This choice makes variable types a reliable contract for pattern parameters and for lowering's dynamic scope: the
composer states intent explicitly, and the compiler does not silently change storage semantics.

\texttt{visit\_var\_decl\_statement} evaluates the initializer, then requires \texttt{is\_assignable(declared\_type,
value\_type)} before registering the symbol.
Assignments to existing variables apply the same rule through the symbol's stored type.

\subsubsection{Expression and Musical Validators}

\noindent \textbf{Binary operators.} Arithmetic operators require numeric operands; comparisons and logical operators
yield \texttt{Bool} via \texttt{binary\_result\_type}.
Adding two sequences remains invalid---sequence concatenation is a structural construct, not \texttt{+}.

\noindent \textbf{\texttt{play} targets.} The play visitor requires musical expression types (\texttt{Note},
\texttt{Chord}, \texttt{Sequence}, or \texttt{Rest}) for melodic sources, and enforces drum-versus-melodic track
consistency when an instrument is declared.
Optional \texttt{:} duration and \texttt{from} offsets must type-check as numeric expressions.

\noindent \textbf{Pattern calls and overload resolution.} When a \texttt{PatternCallExpression} is visited, argument
expressions are typed first.
\texttt{validate\_call} resolves \texttt{find\_pattern\_visible\_by\_signature(callee, argument\_types)}.
Multiple patterns may share a name if their parameter type lists differ; arity alone is insufficient.
If no signature matches, the compiler reports an undefined pattern or a ``no matching overload'' error.
Each argument must satisfy \texttt{is\_assignable(parameter\_type, argument\_type)}.
This replaces version~1.0's arity-only dispatch and removes the need to virtually instantiate pattern bodies in a
second pass.
